% !TEX program = xelatex
% ---------------------------------------------------------------------------
% Dyson 集与伪币识别 —— 十二枚硬币、三次称量，和它背后的编码
% 编译：xelatex Dyson集与伪币识别.tex （跑两遍）
% ---------------------------------------------------------------------------
\documentclass[UTF8,11pt,fontset=macnew]{ctexart}

\usepackage[a4paper,margin=2.4cm]{geometry}
\usepackage{amsmath,amssymb}
\usepackage{booktabs}
\usepackage{array}
\usepackage{enumitem}
\usepackage{xcolor}
\usepackage{tikz}
\usetikzlibrary{arrows.meta,positioning,calc}
\usepackage[most]{tcolorbox}
\usepackage{hyperref}

\definecolor{caramel}{HTML}{A9713C}
\definecolor{rosec}{HTML}{D8678F}
\definecolor{paper}{HTML}{FBF4EA}
\definecolor{inkc}{HTML}{3B2E24}

\hypersetup{colorlinks=true,linkcolor=caramel,urlcolor=rosec,
            pdftitle={Dyson 集与伪币识别},pdfauthor={米露}}

\linespread{1.13}
\setlength{\parskip}{0.35em}
\setlist{itemsep=0.15em,topsep=0.3em}

\newtcolorbox{keybox}[1][小结]{breakable,enhanced,colback=paper,colframe=caramel,
  boxrule=0.7pt,arc=1.5mm,left=3mm,right=3mm,top=2mm,bottom=2mm,
  title={\bfseries #1},coltitle=white,colbacktitle=caramel,fonttitle=\bfseries}

\newtcolorbox{warnbox}[1][注意]{breakable,enhanced,colback=rosec!6,colframe=rosec,
  boxrule=0.7pt,arc=1.5mm,left=3mm,right=3mm,top=2mm,bottom=2mm,
  title={\bfseries #1},coltitle=white,colbacktitle=rosec,fonttitle=\bfseries}

\newcommand{\vv}{\mathbf{v}}
\newcommand{\rr}{\mathbf{r}}

\title{\bfseries Dyson 集与伪币识别\\[0.35em]
  \large 十二枚硬币、三次称量，和它背后的编码}
\author{米露的读书笔记}
\date{2026 年 10 月}

\begin{document}
\maketitle
\pagestyle{plain}
\thispagestyle{empty}

\begin{keybox}[这篇文章想做的事]
十二枚硬币里有一枚伪币，用天平称三次，把它找出来，还要说出它是重是轻。
这是一道流传了很久的智力题，但它的价值不在"想出来"，
而在于它背后有一个非常干净的数学对象：\textbf{Dyson 集}。

这篇文章想讲清楚三件事：为什么每次称量的方案可以写成一组向量；
为什么十二枚是三次称量的\textbf{极限}（多一枚都不行，而"信息量"明明还够）；
以及 Dyson 在 1946 年给出的那套用三进制标签造出这组向量的办法到底在做什么。
（本文由大肥鱼写作）
\end{keybox}

\section{引子：十二枚硬币}

\begin{center}
\begin{tikzpicture}[scale=1,line width=0.9pt,draw=caramel]
  % 底座与立柱
  \draw (-0.55,0) -- (0.55,0);
  \draw (0,0) -- (0,1.9);
  % 横梁
  \draw (-1.7,1.9) -- (1.7,1.9);
  \fill[fill=caramel] (0,1.9) circle (1.6pt);
  % 左右吊盘
  \draw (-1.7,1.9) -- (-1.7,1.35);
  \draw (-2.1,1.35) -- (-1.3,1.35);
  \draw (1.7,1.9) -- (1.7,1.35);
  \draw (1.3,1.35) -- (2.1,1.35);
  \node[font=\small,text=inkc] at (-1.7,1.05) {左盘};
  \node[font=\small,text=inkc] at (1.7,1.05) {右盘};
\end{tikzpicture}
\end{center}

题目是这样的：

\begin{keybox}[十二枚硬币问题]
桌上有 $12$ 枚外观完全相同的硬币。已知其中\textbf{恰好一枚}是伪币，
伪币的重量与真币不同，但\textbf{不知道是偏重还是偏轻}。
用一台\textbf{没有砝码}的天平，最少称几次，能找出伪币并判断它的轻重？
\end{keybox}

答案是 \textbf{3 次}。称 \emph{几次}的问题有信息论的味道，
而\emph{怎么称}的问题则是一道编码题——这正是本文想讲的。

\section{先算一笔信息账}

天平每次称量有三种可能的结果：左盘沉、右盘沉、平衡。
称 $k$ 次，结果的序列最多有 $3^k$ 种。

另一边，要区分的情况有 $12\times 2=24$ 种（哪一枚是伪币 $\times$ 它是重是轻）。
三次称量给出 $3^3=27$ 种结果，而

\[
27 \;>\; 24 .
\]

所以"信息量够用"。

\begin{warnbox}[但信息量够用，不等于办得到]
如果只看到这一步，很容易顺手推出"那 $13$ 枚也行吧"——
$13\times 2 = 26 \le 27$，看着还有富余。
可是 $13$ 枚\textbf{恰恰办不到}，而且原因非常朴素：
天平是一个物理装置，它要求两边放的枚数一样多。
这个约束会吃掉一枚硬币。第 \ref{sec:why12} 节会把这件事说清楚。
\end{warnbox}

所以真正的问题不是数数，而是：\textbf{怎样把方案安排得让 24 种情况
恰好对应 24 种不同的读数，同时每次称量还两边一样重}。
为了回答它，我们需要把"称量方案"翻译成代数。

\section{把称量方案变成向量：Dyson 集}
\label{sec:dyson}

\subsection{给每枚硬币一个向量}

设一共要称 $k$ 次。把整个方案写成一张 $n\times k$ 的表 $W$，
其中第 $i$ 行第 $j$ 列的元素取

\[
W_{ij}=
\begin{cases}
-1, & \text{第 } j \text{ 次称量时，第 } i \text{ 枚硬币放\textbf{左盘}},\\
\ \ 0, & \text{不上秤},\\
+1, & \text{放\textbf{右盘}}.
\end{cases}
\]

这样，第 $i$ 枚硬币就对应一个 $k$ 维向量
\[
\vv_i = (W_{i1},W_{i2},\dots,W_{ik}) \in \{-1,0,+1\}^k .
\]
它的含义是这枚硬币在 $k$ 次称量里的"位置履历"。

再约定读数的符号：

\begin{keybox}[两套符号，方向一致]
\textbf{放置}：$-1$ = 左盘，$0$ = 不上秤，$+1$ = 右盘。\\
\textbf{读数}：$-1$ = 左盘沉，$0$ = 平衡，$+1$ = 右盘沉。\\
两套符号取同样的方向，是为了让下面的结论写成一行：
\[
\text{伪币偏重} \Longrightarrow \rr=\vv_i,
\qquad
\text{伪币偏轻} \Longrightarrow \rr=-\vv_i,
\]
其中 $\rr$ 是三次称量的读数向量。
\end{keybox}

道理很简单：偏重的硬币放在哪一边，哪一边就沉；偏轻的硬币放在哪一边，
那一边就翘起来——读数正好反过来。

于是"能不能判读"这件事，就变成了一句纯粹的代数：

\begin{keybox}[判读条件]
$2n$ 个向量 $\pm\vv_1,\pm\vv_2,\dots,\pm\vv_n$ 必须\textbf{两两不同}。
\end{keybox}

\subsection{三条要求，三个物理含义}

把上面的话拆开，就是下面三条。它们每一条都对应一个物理事实。

\begin{enumerate}
  \item \textbf{每一列的和为 $0$}：$\displaystyle\sum_{i=1}^{n}W_{ij}=0$，
        对每个 $j=1,\dots,k$。
        
        因为 $-1$ 表示左盘、$+1$ 表示右盘，这一列的和就是
        「左盘枚数 $-$ 右盘枚数」。天平要求两边一样多，所以它必须是 $0$。
  \item \textbf{没有零向量}：$\vv_i \neq (0,\dots,0)$。
        
        零向量意味着这枚硬币 $k$ 次都没上过秤。它要是伪币，
        三次全平衡——既判不出它是重是轻，也和"伪币根本不存在"混淆。
  \item \textbf{任意两个向量既不相等，也不互为相反数}：
        $\vv_i \neq \pm\vv_j$（$i\neq j$）。
        
        否则两枚硬币（或者同一枚的重与轻）会产生同一个读数，无法区分。
\end{enumerate}

这三条凑在一起，就是本文的主角：

\begin{keybox}[Dyson 集]
$\{-1,0,+1\}^k$ 中的一族向量 $\{\vv_1,\dots,\vv_n\}$ 称为一个\textbf{Dyson 集}，如果

\begin{enumerate}[label=(\roman*),leftmargin=2.2em]
  \item 每一维各分量之和为 $0$；
  \item 没有零向量；
  \item 所有向量连同它们的反向量 $\pm\vv_i$ 一共 $2n$ 个，两两不同。
\end{enumerate}
\end{keybox}

\begin{warnbox}[为什么是"连同反向量"而不是"互不相同"]
只要求 $\vv_i$ 互不相同是不够的：如果 $\vv_1=-\vv_2$，
那么"第 1 枚偏重"和"第 2 枚偏轻"会给出同一个读数，仍然判不出来。
另外，第 (ii) 条其实被第 (iii) 条包含了——因为 $\mathbf{0}=-\mathbf{0}$，
零向量的两条读数是同一个。这里单列出来，只是为了让"每枚币都要上秤"
这个物理含义不被淹没。
\end{warnbox}

\subsection{一个最小的例子：三枚硬币，两次称量}

在看十二枚之前，先看 $k=2$ 的情形。取

\[
\vv_1=(-1,\,+1),\qquad \vv_2=(0,\,-1),\qquad \vv_3=(+1,\,0).
\]

验一下三条：两列的和分别是 $(-1)+0+(+1)=0$ 与 $(+1)+(-1)+0=0$；
没有零向量；六个向量 $\pm\vv_i$ 两两不同。所以这是一个 $k=2$、$n=3$ 的 Dyson 集。

按"$-1$ 左、$+1$ 右"翻译成称量方案：

\begin{center}
\begin{tabular}{c c c}
\toprule
& 第 1 次称量 & 第 2 次称量 \\
\midrule
硬币 1 & 左盘 & 右盘 \\
硬币 2 & 不上秤 & 左盘 \\
硬币 3 & 右盘 & 不上秤 \\
\midrule
实际称法 & 1 号 vs 3 号 & 2 号 vs 1 号 \\
\bottomrule
\end{tabular}
\end{center}

两次称量各左各右一枚，符合要求。你可以自己试一遍：
无论哪枚是伪币、无论轻重，两次读数拼起来（左沉 / 平衡 / 右沉）
都能唯一指认一种情况——这正是 $6$ 个向量 $\pm\vv_i$ 两两不同的意思。

\section{为什么上限是 12，而不是 13}
\label{sec:why12}

\subsection{宽松的上界：13}

$k$ 次称量共有 $3^k$ 个读数向量，它们都在 $\{-1,0,+1\}^k$ 里。
其中零向量 $(0,\dots,0)$ \textbf{永远不会出现}——因为没有任何一枚硬币的
$\vv_i$ 是零向量，所以伪币总会在至少一次称量里造成倾斜。

去掉零向量，还剩 $3^k-1$ 个。而它们天生成对出现：
$\rr$ 与 $-\rr$ 描述的是同一枚硬币的重与轻。于是可用的情况最多是

\[
\frac{3^k-1}{2}
\quad\xrightarrow{\ k=3\ } \quad \frac{27-1}{2}=13 .
\]

所以只从"结果要够用"来看，$13$ 枚是够的。

\subsection{奇偶性把 13 砍掉}

$13$ 之所以不行，是因为\textbf{天平要求两边枚数相等}这条物理约束
和第 1 维的奇偶性冲突了。看第 $1$ 维（也就是"第 $1$ 次称量"）：

在 $\{-1,0,+1\}^3$ 里，第 1 维等于 $+1$ 的向量有 $3^2=9$ 个，
等于 $-1$ 的也有 $9$ 个，等于 $0$ 的只有 $3^2-1=8$ 个非零向量。
按 $\pm$ 配对，第 1 维\textbf{非零}的那些"对"恰好有 $9$ 对
（每个"非零对"由一个第 1 维为 $+1$ 的向量和它的反向量组成）。

现在假设真的用了 $13$ 枚硬币。因为总共只有 $13$ 对可用，用 $13$ 枚就意味着
\textbf{每一对都要派上用场}，也就是从每一对里挑一个向量当某枚硬币的 $\vv_i$。
于是这 $13$ 个向量中，恰好有 $9$ 个的第 1 维非零。

但第 1 列的和要为 $0$，就要求这 $9$ 个非零的第 1 维里，
$+1$ 和 $-1$ 各占一半。$9$ 是奇数，办不到。矛盾。

\begin{keybox}[一般结论]
$k$ 次称量能处理的最多硬币数（伪币必存在、轻重未知）是
\[
n_{\max}=\frac{3^k-3}{2},
\]
而"信息量"给出的上界是 $(3^k-1)/2$。两者恰好差一枚，
差的这一枚就是被"两边要一样重"这条物理约束挤掉的。
\end{keybox}

\begin{center}
\begin{tabular}{c c c c}
\toprule
称量次数 $k$ & 信息量上界 $\frac{3^k-1}{2}$ & 真正可达 $\frac{3^k-3}{2}$ & 每次称量两边各放 \\
\midrule
2 & 4 & 3   & 1 枚 \\
3 & 13 & \textbf{12} & 4 枚 \\
4 & 40 & 39  & 13 枚 \\
5 & 121 & 120 & 40 枚 \\
\bottomrule
\end{tabular}
\end{center}

（"两边各放"那一列就是 $n_{\max}/3$，下一节会解释它为什么一定是 $3$ 的倍数。）
上界只是上界，还得真的造得出来——这正是 Dyson 在 1946 年做的事。

\section{Dyson 的构造：用三进制标签给硬币编号}
\label{sec:construct}

Dyson 的办法出人意料地简单：\textbf{别去凑向量，去数数。}

\subsection{第一步：把所有 n 位三进制数摆出来}

设要称 $n$ 次\footnote{这里沿用 Dyson 原文把称量次数记作 $n$；
本文前面用的是 $k$，两者是同一个东西。}，目标是处理
$M=\dfrac{3^n-3}{2}$ 枚硬币。

把 $0$ 到 $3^n-1$ 全都写成 $n$ 位三进制（不足的前面补 $0$），
得到 $3^n$ 个"标签"。例如 $n=3$ 时，共有 $27$ 个标签。

\subsection{第二步：互补配对}

定义一个标签的\textbf{互补标签}：把它的 $0$ 全换成 $2$、$2$ 全换成 $0$，
$1$ 保持不动。（也就是用 $3^n-1$ 去减它。）

一个标签和它的互补标签配成一对，这一对就是\textbf{一枚硬币}。
为什么？因为互补在下面的翻译里恰好对应取反：
数字 $0,1,2$ 将分别翻译成 $-1,0,+1$，而 $0\leftrightarrow 2$ 正是 $-1\leftrightarrow+1$。
所以"一对互补标签"就是"一个 $\pm\vv$"，也就是一枚硬币的重与轻两种读法。

\subsection{第三步：抛掉三个恒定标签}

不能把 $3^n$ 个标签全用上，因为\textbf{互补配对配不齐}。
问题出在三个"每一位都相同"的标签上：

\begin{itemize}
  \item $\underbrace{11\cdots1}_{n}$：它的互补标签是它自己，配不成对，只能丢掉。
  \item $\underbrace{00\cdots0}_{n}$ 与 $\underbrace{22\cdots2}_{n}$：它们互为互补，
        看似可以配成一对，但它们是"永远放左盘"和"永远放右盘"。
        更直接的问题是：它们没有任何"数字变化"，
        下一步要用的顺时针判定对它们失效（见下）。
\end{itemize}

丢掉这 $3$ 个，剩下 $3^n-3$ 个标签，恰好两两配成

\[
\frac{3^n-3}{2}=M
\]

对，也就是 $M$ 枚硬币。这就是 $M$ 的来历。

\subsection{第四步：顺时针与逆时针——从每对里挑一个}

现在每枚硬币有两个标签（互为互补）。为了说话方便，需要\textbf{统一从每对里挑一个}。
Dyson 的挑法是：

\begin{keybox}[顺时针 / 逆时针]
从左往右读一个标签，找到\textbf{第一个与首位不同的数字}，看这次变化的方向：
\[
0\to1,\quad 1\to2,\quad 2\to0
\]
叫\textbf{顺时针}（数字在 $0\to1\to2\to0$ 这个循环上前进了一格）；
反过来的 $1\to0$、$2\to1$、$0\to2$ 叫\textbf{逆时针}。
\end{keybox}

比如标签 $012$：首位是 $0$，第一个不同的是第 2 位的 $1$，变化是 $0\to1$，
所以它是顺时针。标签 $210$：首位 $2$，第 2 位是 $1$，变化 $2\to1$，
是逆时针。又比如 $120$：首位 $1$，第 2 位 $2$，变化 $1\to2$，
顺时针。

关键性质是：\textbf{一对互补标签里，恰好一个是顺时针，一个是逆时针}，
所以"取顺时针的那个"永远能唯一地选出代表，不会打架。
（至于那三个恒定标签——它们没有任何变化，压根判不出方向，
所以在这一步被正式排除。）

\subsection{第五步：把标签翻译成称量}

规则是：把三进制数字翻译成放置方式，
\[
0 \longmapsto -1\ (\text{左盘}), \qquad
1 \longmapsto 0\ (\text{不上秤}), \qquad
2 \longmapsto +1\ (\text{右盘}).
\]

于是，第 $i$ 次称量就是：

\begin{keybox}[Dyson 的称量方案]
在每枚硬币的\textbf{顺时针标签}里看第 $i$ 位的数字：\\
数字为 $0$ 的硬币放\textbf{左盘}；数字为 $2$ 的放\textbf{右盘}；
数字为 $1$ 的\textbf{不上秤}。
\end{keybox}

\subsection{为什么天平真的会平衡}

这是整个构造里最漂亮的一步。它靠的是"顺时针"这个判定在
数字循环置换下不变。

设 $\sigma$ 表示三进制数字的循环置换 $0\to1\to2\to0$。
把一枚硬币的顺时针标签的每一位都做 $\sigma$，得到的还是顺时针标签——
因为 $\sigma$ 把三种"前进"的变化一一对应地变成了"前进"的变化：
$0\to1$ 变成 $1\to2$，$1\to2$ 变成 $2\to0$，$2\to0$ 变成 $0\to1$。
反过来，"后退"的变化也照样是后退。

所以 $\sigma$ 把全体顺时针标签映到全体顺时针标签，而且是双射。
于是它把"顺时针标签第 $i$ 位是 $0$ 的硬币"整体搬到了
"顺时针标签第 $i$ 位是 $1$ 的硬币"，再搬到"是 $2$"的那一批，最后回到原处。
三批硬币因此\textbf{一样多}，各有
\[
\frac{M}{3}=\frac{3^{n-1}-1}{2}
\]
枚。第 $i$ 次称量时，左盘放的是第 $i$ 位为 $0$ 的那一批，
右盘放的是第 $i$ 位为 $2$ 的那一批，两者枚数相同——天平正好放平。

\begin{keybox}[回头看：构造出的向量确实是 Dyson 集]
把顺时针标签翻译成 $\{-1,0,+1\}^n$ 里的向量，逐条对一下三条要求：
列和为 $0$ 就是刚才证的"左右一样多"；
没有零向量是因为零向量要求标签全是 $1$，而 $\underbrace{11\cdots1}_{n}$ 已经被丢掉了；
两两不相等、也不互为相反数，则是因为每个标签在全体标签中只出现一次，
而"互为相反数"恰好就是"互为互补标签"，那一对只给了一枚硬币。
\end{keybox}

\section{十二枚硬币的具体方案}

把 $n=3$ 代进去：$M=(27-3)/2=12$，每次称量两边各 $4$ 枚，
一共三次。下面把整套东西列出来。

\subsection{标签表}

\begin{center}
\begin{tabular}{c c c c}
\toprule
硬币 & 顺时针标签 & 互补标签 & 向量 $\vv_i$ \\
\midrule
1  & 001 & 221 & $(-1,-1,\ \ 0)$ \\
2  & 220 & 002 & $(\ \ 1,\ \ 1,-1)$ \\
3  & 010 & 212 & $(-1,\ \ 0,-1)$ \\
4  & 011 & 211 & $(-1,\ \ 0,\ \ 0)$ \\
5  & 012 & 210 & $(-1,\ \ 0,\ \ 1)$ \\
6  & 202 & 020 & $(\ \ 1,-1,\ \ 1)$ \\
7  & 201 & 021 & $(\ \ 1,-1,\ \ 0)$ \\
8  & 200 & 022 & $(\ \ 1,-1,-1)$ \\
9  & 122 & 100 & $(\ \ 0,\ \ 1,\ \ 1)$ \\
10 & 121 & 101 & $(\ \ 0,\ \ 1,\ \ 0)$ \\
11 & 120 & 102 & $(\ \ 0,\ \ 1,-1)$ \\
12 & 112 & 110 & $(\ \ 0,\ \ 0,\ \ 1)$ \\
\bottomrule
\end{tabular}
\end{center}

这 $12$ 个向量就是一组 $n=3$ 的 Dyson 集：三列的和都是 $0$，
没有零向量，$24$ 个 $\pm\vv_i$ 两两不同。
（注意 $13$ 对里只用了 $12$ 对，剩下没用到的那一对是
$(-1,-1,-1)$ 与 $(1,1,1)$——正是被丢掉的那两个恒定标签 $000$ 和 $222$。）

\subsection{三次称量}

对着标签表读第 $i$ 位，$0$ 放左、$2$ 放右、$1$ 不上秤：

\begin{center}
\begin{tabular}{c c c c}
\toprule
& 左盘 & 右盘 & 不上秤 \\
\midrule
第 1 次 & 1, 3, 4, 5   & 2, 6, 7, 8    & 9, 10, 11, 12 \\
第 2 次 & 1, 6, 7, 8   & 2, 9, 10, 11  & 3, 4, 5, 12 \\
第 3 次 & 2, 3, 8, 11  & 5, 6, 9, 12   & 1, 4, 7, 10 \\
\bottomrule
\end{tabular}
\end{center}

每次左右各 $4$ 枚。

\begin{warnbox}[这个方案是"非自适应"的]
三次称量\textbf{事先全部定好}，不需要根据前一次的结果临时决定下一次怎么称。
这比题目要求的（允许自适应）更强。有些解法会让第二次称量依赖于第一次的结果，
那样做当然也可以，但 Dyson 集给出的是这种"一次性排好"的方案。
\end{warnbox}

\subsection{怎么判读}

称完之后，把三次结果按"左沉 / 平衡 / 右沉"记成 $-1/0/+1$，
拼成一个读数向量 $\rr$。然后：

\begin{keybox}[判读规则]
在上面的向量表里查 $\rr$：
\begin{itemize}
  \item 若 $\rr$ 正好等于某个 $\vv_i$，则第 $i$ 枚是伪币，且\textbf{偏重}；
  \item 若 $\rr$ 正好等于某个 $-\vv_i$，则第 $i$ 枚是伪币，且\textbf{偏轻}。
\end{itemize}
两种情况必居其一，因为可用的 $24$ 个读数向量都用上了。
\end{keybox}

举个例子：称出来是「左沉、平衡、右沉」，即 $\rr=(-1,0,+1)$。
查表，$\vv_5=(-1,0,+1)$，所以\textbf{第 5 枚偏重}。
如果是「右沉、平衡、左沉」，即 $\rr=(1,0,-1)=-\vv_5$，
那就是\textbf{第 5 枚偏轻}。

完整的对照表如下。读数按「第 1 次、第 2 次、第 3 次」的顺序写成三个符号
（$- $ 左沉，$0$ 平衡，$+$ 右沉），每一行是一枚硬币的两种可能读法：

\begin{center}
\footnotesize
\setlength{\tabcolsep}{6pt}
\begin{tabular}{c l @{\hspace{2.4em}} c l}
\toprule
读数（$\vv_i$） & 结论 & 读数（$-\vv_i$） & 结论 \\
\midrule
${-}{-}0$ & 1 号偏重  & ${+}{+}0$ & 1 号偏轻 \\
${+}{+}{-}$ & 2 号偏重 & ${-}{-}{+}$ & 2 号偏轻 \\
${-}0{-}$ & 3 号偏重  & ${+}0{+}$ & 3 号偏轻 \\
${-}00$ & 4 号偏重    & ${+}00$ & 4 号偏轻 \\
${-}0{+}$ & 5 号偏重  & ${+}0{-}$ & 5 号偏轻 \\
${+}{-}{+}$ & 6 号偏重 & ${-}{+}{-}$ & 6 号偏轻 \\
${+}{-}0$ & 7 号偏重  & ${-}{+}0$ & 7 号偏轻 \\
${+}{-}{-}$ & 8 号偏重 & ${-}{+}{+}$ & 8 号偏轻 \\
$0{+}{+}$ & 9 号偏重  & $0{-}{-}$ & 9 号偏轻 \\
$0{+}0$ & 10 号偏重   & $0{-}0$ & 10 号偏轻 \\
$0{+}{-}$ & 11 号偏重 & $0{-}{+}$ & 11 号偏轻 \\
$00{+}$ & 12 号偏重   & $00{-}$ & 12 号偏轻 \\
\bottomrule
\end{tabular}
\end{center}

顺带一提：$(\pm1,\pm1,\pm1)$ 型的读数只有 $6$ 个会出现，
$(-1,-1,-1)$ 和 $(1,1,1)$ 永远不会出现——它们正是被丢掉的那一对恒定标签。

\section{小结}

\subsection{这套东西的骨架}

\begin{enumerate}
  \item 把称量方案写成 $n\times k$ 的 $\{-1,0,+1\}$ 矩阵，
        每一行是一枚硬币的向量 $\vv_i$。
  \item 三条要求——列和为 $0$、无零向量、$\pm\vv_i$ 两两不同——
        合起来就是"$\{\vv_i\}$ 是一个 Dyson 集"。
  \item 计数：$3^k$ 个读数去掉零向量，按 $\pm$ 配对，最多 $(3^k-1)/2$。
  \item 奇偶性：用满 $(3^k-1)/2$ 时，第 1 维非零的向量有 $3^{k-1}$ 个，
        $k\ge2$ 时是奇数，无法让天平两边相等。于是上界降到 $(3^k-3)/2$。
  \item Dyson 的构造证明这个上界\textbf{可以达到}：用全部 $n$ 位三进制标签，
        丢掉 $000\cdots0$、$111\cdots1$、$222\cdots2$，互补配对，
        每对取顺时针的那一个。
\end{enumerate}

\subsection{几点延伸}

\textbf{为什么"顺时针"这一步不能省。}
它看起来只是一个技术性的约定，实际上承担了两件事：
一是从每对互补标签里\textbf{一致地}挑出一个代表（否则同一枚币在不同称量里
可能被取成相反的方向，列和就乱了）；
二是让"循环置换 $\sigma$"有落脚点，从而一句话证明每次称量两边一样多。
Dyson 当年用这个巧劲，避开了繁琐的构造性讨论。

\textbf{Dyson 集不是唯一的。}
本文给出的 $12$ 个向量只是其中一个选择。换一套代表（比如把某几枚硬币的
向量同时取反再调换称量的左右），只要三条要求还成立，就还是合法的方案。
文献里常见的那组是
\[
\begin{aligned}
\{&(-1,-1,0),\ (0,0,1),\ (1,1,-1),\ (-1,0,-1),\ (0,1,0),\ (1,-1,1),\\
 &(0,-1,-1),\ (1,0,0),\ (-1,1,1),\ (-1,0,1),\ (0,1,-1),\ (1,-1,0)\},
\end{aligned}
\]
和本文的这组用满了同样多的"对"，只是换掉了没用到的那一对。

\textbf{"信息量够"与"可达"之间的那道缝。}
$(3^k-1)/2$ 和 $(3^k-3)/2$ 只差一枚，但这一枚的差距是真实的，
而且原因（两边枚数必须相等 $\Rightarrow$ 某一位非零的分量个数必须是偶数）
和香农式信息论毫无关系。这类"信息论下界不紧"的现象在编码与组合设计里很常见，
也是这道老题真正有意思的地方。

\begin{keybox}[参考]
\begin{itemize}
  \item F.~J.~Dyson, \textit{The Problem of the Pennies},
        The Mathematical Gazette \textbf{30} (1946), 231--234. \\
        三进制标签与顺时针/逆时针的原始构造。
  \item T.~Tokuda, Y.~Watanabe,
        \textit{A bounded partition approach to identifying one fake coin and its type},
        arXiv:2306.10733 (2023). \\
        用现代语言重述 Dyson 集的定义，并给出不依赖 Dyson 集的递归解法。
\end{itemize}
\end{keybox}

\end{document}
